\documentclass[12pt, a4paper]{article}
\usepackage[italian]{babel}
\usepackage[utf8]{inputenc}
\usepackage[T1]{fontenc}
\usepackage{amsmath}
\usepackage{amssymb}
\usepackage{algorithm}
\usepackage{algpseudocode}
\usepackage{listings}
\usepackage{xcolor}
\usepackage{graphicx}
\usepackage{booktabs}
\usepackage[hidelinks]{hyperref}
\usepackage{geometry}
\usepackage{float}
\usepackage{caption}

\geometry{top=3cm, bottom=3cm, left=3cm, right=3cm}

\definecolor{codegreen}{rgb}{0,0.6,0}
\definecolor{codegray}{rgb}{0.5,0.5,0.5}
\definecolor{codepurple}{rgb}{0.58,0,0.82}
\definecolor{backcolour}{rgb}{0.97,0.97,0.97}

\lstdefinestyle{pythonstyle}{
    backgroundcolor=\color{backcolour},
    commentstyle=\color{codegreen}\itshape,
    keywordstyle=\color{blue}\bfseries,
    numberstyle=\tiny\color{codegray},
    stringstyle=\color{codepurple},
    basicstyle=\ttfamily\small,
    breakatwhitespace=false,
    breaklines=true,
    captionpos=b,
    keepspaces=true,
    numbers=left,
    numbersep=5pt,
    showspaces=false,
    showstringspaces=false,
    showtabs=false,
    tabsize=4,
    language=Python,
    frame=single,
    rulecolor=\color{codegray}
}
\lstset{style=pythonstyle}
\renewcommand{\lstlistingname}{Listato}

\title{\textbf{Retrieval di Immagini tramite Vector Quantization}\\[0.5em]
\large Implementazione e Confronto tra Flat VQ e Tree-Structured VQ\\[0.3em]
sul Dataset Corel-1K}
\author{Francesco Paolo Sferratore}
\date{16/03/2026}

\begin{document}

\maketitle
\tableofcontents
\newpage

% ==============================================================================
\section{Introduzione}
% ==============================================================================

Il \textit{recupero di immagini basato sul contenuto} (Content-Based Image Retrieval, CBIR)
studia come trovare automaticamente immagini simili in un archivio, senza fare affidamento
su descrizioni testuali o tag assegnati manualmente. Dato un archivio e un'immagine di
interrogazione (\textit{query}), l'obiettivo è restituire le immagini dell'archivio ordinate
per somiglianza visiva.

Il presente lavoro implementa un sistema CBIR basato sulla \textit{Vector Quantization} (VQ).
L'idea è descrivere ogni immagine con un insieme compatto di vettori rappresentativi
(\textit{codeword}), detti \textit{codebook}. Due immagini con colori simili produrranno
codebook simili; confrontando il codebook di un'immagine del database con i vettori di
colore della query si ottiene una misura di quanto le due immagini si assomigliano.

\subsection*{Dataset}
Il sistema è stato valutato sul dataset \textbf{Corel-1K}, un benchmark standard per CBIR.
Contiene 1000 immagini suddivise in 10 categorie semantiche, 100 immagini per categoria:
\begin{center}
\textit{africans, beaches, buildings, buses, dinosaurs, elephants, flowers, food, horses, mountains}
\end{center}
Le immagini rilevanti per una query sono quelle della stessa categoria semantica.

\subsection*{Struttura del lavoro}
Sono state implementate e confrontate due varianti:
\begin{itemize}
    \item \textbf{Flat VQ}: codebook piatto con $K=8$ centroidi per regione, costruito
          tramite k-means.
    \item \textbf{TSVQ} (Tree-Structured VQ): codebook a struttura ad albero binario di
          profondità $d=3$, che produce $2^3 = 8$ codeword foglia per regione.
\end{itemize}
Entrambi producono 8 codeword per regione, rendendo il confronto equo.

Il lavoro è organizzato come segue: la Sezione~\ref{sec:flatvq} descrive l'algoritmo Flat VQ,
la Sezione~\ref{sec:impl_flatvq} la sua implementazione, la Sezione~\ref{sec:tsvq} l'algoritmo
TSVQ, la Sezione~\ref{sec:impl_tsvq} la sua implementazione, la Sezione~\ref{sec:valutazione}
il confronto teorico, e la Sezione~\ref{sec:evaluation} i risultati sperimentali.

% ==============================================================================
\section{Algoritmo Flat VQ}
\label{sec:flatvq}
% ==============================================================================

\subsection{Spazio colore CIE L*u*v*}

Le immagini vengono convertite da RGB allo spazio colore \textbf{CIE L*u*v*} prima di
qualsiasi elaborazione. CIE L*u*v* è uno spazio percettivamente uniforme: distanze uguali
in questo spazio corrispondono a differenze di colore percepite come uguali dall'occhio
umano. In RGB questo non vale: due coppie di colori con la stessa distanza numerica possono
sembrare all'occhio molto diverse tra loro.

Il canale $L^*$ rappresenta la luminosità, mentre $u^*$ e $v^*$ rappresentano le componenti
cromatiche.

\subsection{Suddivisione in blocchi}

Ogni immagine viene ridimensionata a \textbf{256$\times$256 pixel} per uniformare la
risoluzione di tutte le immagini nel database. L'immagine viene poi divisa in blocchi
non sovrapposti di \textbf{4$\times$4 pixel}. Con una risoluzione di 256$\times$256 si
ottengono $64 \times 64 = 4096$ blocchi per immagine, ciascuno contenente 16 pixel.

La dimensione 4$\times$4 è un compromesso: blocchi più piccoli aumentano la risoluzione
spaziale ma rendono instabili le stime statistiche (troppo pochi pixel); blocchi più grandi
stabilizzano le stime ma perdono il dettaglio locale.

\subsection{Vettore di feature per blocco}

Per ogni blocco si calcolano la \textbf{media} e la \textbf{varianza} del colore per
ciascuno dei tre canali L*, u*, v*. Si ottiene un vettore di 6 valori:
\[
\mathbf{f} = [\,\mu_L,\; \mu_u,\; \mu_v,\; \sigma^2_L,\; \sigma^2_u,\; \sigma^2_v\,]
\]
La \textbf{media} descrive il colore dominante del blocco. La \textbf{varianza} misura
quanto il colore varia all'interno del blocco: un valore basso indica un blocco uniforme
(ad esempio un pezzo di cielo), un valore alto indica variazione cromatica (ad esempio
un contorno o una texture).

Si è scelto di non includere la \textbf{skewness} (il terzo momento statistico), nonostante
possa catturare asimmetrie nella distribuzione del colore. Su blocchi 4$\times$4 di soli
16 pixel la sua stima risulta troppo rumorosa per essere discriminativa; i test sperimentali
hanno confermato che includerla peggiorava i risultati.

\subsection{Normalizzazione delle feature}

Media e varianza operano su scale molto diverse: i valori medi appartengono all'intervallo
$[0, 255]$, mentre le varianze possono arrivare fino a $255^2/4 = 16256$ nel caso peggiore.
Senza correzione, le varianze dominerebbero completamente il calcolo delle distanze.

Per rendere tutte le dimensioni comparabili, ogni feature viene divisa per il suo massimo
teorico:
\begin{itemize}
    \item medie: divise per 255
    \item varianze: divise per 16256
\end{itemize}
Dopo la normalizzazione tutti i valori appartengono all'intervallo $[0, 1]$ e contribuiscono
equamente alle distanze.

\subsection{Griglia spaziale}

I 4096 blocchi dell'immagine vengono assegnati a una \textbf{griglia 4$\times$4 di regioni},
per un totale di 16 regioni. Ogni regione contiene $16 \times 16 = 256$ blocchi.

La griglia spaziale codifica la posizione: senza di essa, due immagini con gli stessi colori
ma disposti diversamente nello spazio sarebbero indistinguibili. Con la griglia, il cielo in
alto e il prato in basso contribuiscono a regioni diverse, permettendo al sistema di
distinguere composizioni spazialmente differenti. Le regioni sono numerate in ordine
riga-per-riga da 0 a 15.

\subsection{Costruzione del codebook con k-means}

Per ogni regione di ogni immagine del database, si applica l'algoritmo \textbf{k-means} con
$K=8$ cluster ai 256 vettori di feature di quella regione. K-means trova 8 centroidi che
rappresentano i gruppi di colori più ricorrenti nella regione; l'insieme di questi 8 centroidi
è il \textit{codebook} della regione.

Poiché k-means è sensibile al punto di partenza casuale, viene eseguito \textbf{5 volte}
(\texttt{n\_init=5}) mantenendo il risultato con distorsione minima. Ogni immagine è
descritta da $16 \times 8 = 128$ centroidi in totale.

\subsection{Misura di similarità: MSE di quantizzazione}

Il principio del retrieval è il seguente: se un'immagine del database $I$ è simile alla
query $Q$, il codebook di $I$ -- costruito per descrivere i colori di $I$ -- riuscirà a
rappresentare bene anche i colori di $Q$, producendo un errore di approssimazione basso.

Per misurare questo errore, si prende ogni blocco della query e si trova il codeword del
database più vicino. La distanza al quadrato tra il blocco e il suo codeword più vicino è
l'errore di quantizzazione di quel blocco. Il \textbf{punteggio di dissimilarità} tra query
e immagine è la media di questi errori su tutti i blocchi e tutte le regioni: più è basso,
più le due immagini si assomigliano. Le immagini del database vengono ordinate per punteggio
crescente e le prime vengono restituite come risultato.

% ==============================================================================
\section{Implementazione Flat VQ}
\label{sec:impl_flatvq}
% ==============================================================================

Il sistema Flat VQ è implementato in quattro script Python: \texttt{preprocess.py},
\texttt{build\_codebooks.py}, \texttt{retrieve.py} ed \texttt{evaluate.py}.

\subsection{Modulo di preprocessing: \texttt{preprocess.py}}

Ridimensiona ogni immagine grezza a 256$\times$256 pixel e la salva in
\texttt{data/processed/}, preservando la struttura delle sottocartelle di categoria.

\begin{lstlisting}[caption={Ciclo principale di \texttt{preprocess.py}}]
for img_path in input_dir.rglob("*"):
    if img_path.suffix.lower() not in [".jpg", ".jpeg", ".png"]:
        continue
    img = cv2.imread(str(img_path))
    img = cv2.resize(img, (size, size))
    save_path = save_dir / (img_path.stem + ".jpg")
    cv2.imwrite(str(save_path), img)
\end{lstlisting}

\textbf{Scelte implementative:}
\begin{itemize}
    \item Le immagini vengono salvate in formato BGR senza convertirle in L*u*v*. La
          conversione di spazio colore avviene in memoria in \texttt{build\_codebooks.py},
          appena prima dell'estrazione delle feature, evitando di degradare le immagini
          salvate con una ricompressione in uno spazio non standard.
    \item Il ridimensionamento usa l'interpolazione bilineare predefinita di OpenCV,
          adeguata per questo tipo di applicazione.
\end{itemize}

\subsection{Modulo di costruzione dei codebook: \texttt{build\_codebooks.py}}

Questo script costruisce i codebook per tutte le immagini del database e li salva su disco.

\subsubsection{Parametri globali}

\begin{lstlisting}[caption={Parametri di configurazione di \texttt{build\_codebooks.py}}]
BLOCK_SIZE = 4
GRID       = 4
K          = 8
NORM_SCALE = np.array([255, 255, 255, 16256, 16256, 16256], dtype=np.float64)
\end{lstlisting}

I parametri sono definiti come costanti globali. Devono essere identici in
\texttt{retrieve.py}: codebook costruiti con certi parametri sono utilizzabili solo con
lo stesso preprocessing nella fase di query.

\subsubsection{Funzione \texttt{extract\_block\_features}}

\begin{lstlisting}[caption={Funzione \texttt{extract\_block\_features}}]
def extract_block_features(img):
    h, w, _ = img.shape
    feats = []
    for y in range(0, h, BLOCK_SIZE):
        for x in range(0, w, BLOCK_SIZE):
            block = img[y:y+BLOCK_SIZE, x:x+BLOCK_SIZE]
            mean = block.mean(axis=(0, 1))
            var  = block.var(axis=(0, 1))
            feats.append(np.concatenate([mean, var]))
    feats = np.array(feats)
    return feats / NORM_SCALE
\end{lstlisting}

La funzione scorre l'immagine blocco per blocco con due cicli annidati. Per ogni blocco
4$\times$4 calcola media e varianza per canale (L*, u*, v*), concatena i sei valori in un
vettore e lo aggiunge alla lista. Al termine, divide l'intero array per \texttt{NORM\_SCALE}
in un'unica operazione vettoriale, producendo un array di forma $(4096, 6)$.

\subsubsection{Funzione \texttt{assign\_regions}}

\begin{lstlisting}[caption={Funzione \texttt{assign\_regions}}]
def assign_regions(n_blocks_y, n_blocks_x):
    region_map = np.zeros((n_blocks_y, n_blocks_x), dtype=int)
    ry = n_blocks_y // GRID
    rx = n_blocks_x // GRID
    r = 0
    for gy in range(GRID):
        for gx in range(GRID):
            y0, y1 = gy * ry, (gy + 1) * ry
            x0, x1 = gx * rx, (gx + 1) * rx
            region_map[y0:y1, x0:x1] = r
            r += 1
    return region_map
\end{lstlisting}

Costruisce una mappa 2D della stessa forma della griglia di blocchi, dove ogni cella contiene
l'ID di regione (da 0 a 15) del blocco corrispondente. L'assegnazione avviene tramite slicing
NumPy sull'area rettangolare di ogni regione, senza loop sui singoli blocchi.

\subsubsection{Funzione \texttt{build\_codebooks}}

\begin{lstlisting}[caption={Funzione \texttt{build\_codebooks} (Flat VQ)}]
def build_codebooks(img):
    h, w, _ = img.shape
    n_blocks_y = h // BLOCK_SIZE
    n_blocks_x = w // BLOCK_SIZE
    features   = extract_block_features(img)
    region_ids = assign_regions(n_blocks_y, n_blocks_x).flatten()
    codebooks = []
    for r in range(GRID * GRID):
        region_feats = features[region_ids == r]
        k = min(K, len(region_feats))
        kmeans = KMeans(n_clusters=k, n_init=5)
        kmeans.fit(region_feats)
        codebooks.append(kmeans.cluster_centers_)
    return codebooks
\end{lstlisting}

La funzione estrae le feature di tutti i blocchi, ne determina la regione di appartenenza,
e per ognuna delle 16 regioni esegue k-means sui soli 256 vettori di quella regione.

La riga \texttt{k = min(K, len(region\_feats))} è una guardia di robustezza: impedisce di
chiamare k-means con più cluster che punti disponibili, cosa che causerebbe un errore.

\subsubsection{Ciclo principale e salvataggio}

\begin{lstlisting}[caption={Ciclo principale di \texttt{build\_codebooks.py}}]
for file in INPUT_DIR.rglob("*.jpg"):
    img = cv2.imread(str(file))
    img = cv2.cvtColor(img, cv2.COLOR_BGR2Luv)
    codebooks = build_codebooks(img)
    relative = file.relative_to(INPUT_DIR)
    save_dir  = OUTPUT_DIR / relative.parent
    save_dir.mkdir(parents=True, exist_ok=True)
    save_path = save_dir / (file.stem + ".npz")
    np.savez(save_path, codebooks=codebooks)
\end{lstlisting}

Per ogni immagine, la conversione in L*u*v* avviene in memoria appena prima dell'estrazione
delle feature. Il percorso di salvataggio replica la struttura delle sottocartelle originali,
in modo che la categoria di ogni immagine sia leggibile direttamente dal percorso del file
\texttt{.npz}, senza bisogno di file di metadati separati.

\subsection{Modulo di recupero: \texttt{retrieve.py}}

Implementa la pipeline di query: data un'immagine di interrogazione, restituisce la lista
delle immagini del database ordinate per similarità.

\subsubsection{Funzioni \texttt{preprocess} ed \texttt{extract\_block\_features}}

Queste funzioni sono \textbf{identiche} alle corrispondenti in \texttt{build\_codebooks.py}.
È un requisito fondamentale: se i vettori di feature della query fossero estratti in modo
diverso da quelli del database, le distanze calcolate non avrebbero senso.

\subsubsection{Funzione \texttt{get\_query\_signature}}

\begin{lstlisting}[caption={Funzione \texttt{get\_query\_signature}}]
def get_query_signature(img):
    h, w, _ = img.shape
    n_blocks_y = h // BLOCK_SIZE
    n_blocks_x = w // BLOCK_SIZE
    features   = extract_block_features(img)
    region_ids = assign_regions(n_blocks_y, n_blocks_x).flatten()
    signature = []
    for r in range(GRID * GRID):
        region_feats = features[region_ids == r]
        signature.append(region_feats)
    return signature
\end{lstlisting}

Per la query \textbf{non viene eseguito k-means}. Si raccolgono semplicemente i vettori di
feature raggruppati per regione. Questi vettori verranno poi confrontati con i codeword dei
codebook del database durante la fase di scoring. La funzione restituisce una lista di 16
array, uno per regione, ciascuno di forma $(256, 6)$.

\subsubsection{Funzione \texttt{region\_mse}}

\begin{lstlisting}[caption={Funzione \texttt{region\_mse} (Flat VQ) con broadcasting NumPy}]
def region_mse(query_vectors, codebook):
    dists = np.sum(
        (query_vectors[:, np.newaxis, :] - codebook) ** 2,
        axis=2
    )  # shape: (n, K)
    return float(np.mean(dists.min(axis=1)))
\end{lstlisting}

Per ogni blocco della query, trova il codeword del database più vicino e registra la distanza
al quadrato. Il risultato finale è la media di queste distanze su tutti i blocchi della
regione.

\textbf{Ottimizzazione vettoriale}: invece di iterare in Python su ciascuno dei 256 vettori
uno alla volta, tutti i confronti (256 vettori contro 8 codeword) vengono eseguiti in
un'unica operazione NumPy tramite broadcasting. Questo riduce le chiamate al motore NumPy
da 256 a 1, eliminando l'overhead di invocazione ripetuta su array piccoli. Considerando
che questa funzione viene chiamata 16 milioni di volte in una valutazione completa, il
risparmio è decisivo: il tempo di esecuzione scende da giorni a meno di un'ora.

\subsubsection{Funzione \texttt{score\_against\_db\_image}}

\begin{lstlisting}[caption={Funzione \texttt{score\_against\_db\_image}}]
def score_against_db_image(query_signature, codebooks):
    regional_mses = []
    for r in range(GRID * GRID):
        mse = region_mse(query_signature[r], codebooks[r])
        regional_mses.append(mse)
    return float(np.mean(regional_mses))
\end{lstlisting}

Calcola l'MSE per ognuna delle 16 regioni e restituisce la loro media come punteggio
complessivo di dissimilarità tra la query e un'immagine del database.

\subsubsection{Funzione \texttt{load\_all\_codebooks}}

\begin{lstlisting}[caption={Funzione \texttt{load\_all\_codebooks}}]
def load_all_codebooks():
    db = []
    for cb_file in CODEBOOKS_DIR.rglob("*.npz"):
        data = np.load(cb_file, allow_pickle=True)
        db.append({
            'codebooks': data['codebooks'],
            'stem':      cb_file.stem,
            'category':  cb_file.parent.name,
        })
    return db
\end{lstlisting}

Carica in RAM tutti i codebook del database. Viene chiamata una sola volta prima del ciclo
di query: senza questa ottimizzazione, ogni query rileggeva da disco gli stessi 1000 file,
per un totale di un milione di letture ridondanti in una valutazione completa.

\subsubsection{Funzione \texttt{retrieve}}

\begin{lstlisting}[caption={Funzione \texttt{retrieve}}]
def retrieve(query_path, top_k=10, db=None):
    img             = preprocess(query_path)
    query_signature = get_query_signature(img)
    if db is None:
        db = load_all_codebooks()
    results = []
    for entry in db:
        score = score_against_db_image(query_signature, entry['codebooks'])
        results.append({
            'image':    entry['stem'],
            'category': entry['category'],
            'score':    score,
        })
    results.sort(key=lambda x: x['score'])
    for i, r in enumerate(results, 1):
        r['rank'] = i
    return results[:top_k]
\end{lstlisting}

Esegue la pipeline completa: preprocessa la query, ne estrae la firma, calcola il punteggio
contro ogni immagine del database, ordina per punteggio crescente e restituisce le prime
\texttt{top\_k}. Il parametro \texttt{db} consente di passare i codebook già caricati in
memoria, oppure di ometterlo per un utilizzo autonomo su singola query.

\subsection{Modulo di valutazione: \texttt{evaluate.py}}

Misura la qualità del retrieval sull'intero dataset, usando ogni immagine a turno come query
e valutando le altre 999 come candidati.

\subsubsection{Funzione \texttt{compute\_pr\_curve}}

\begin{lstlisting}[caption={Funzione \texttt{compute\_pr\_curve}}]
def compute_pr_curve(results, query_stem, query_category):
    precisions, recalls = [], []
    n_correct, rank = 0, 0
    for r in results:
        if r['image'] == query_stem:
            continue
        rank += 1
        if r['category'] == query_category:
            n_correct += 1
        precisions.append(n_correct / rank)
        recalls.append(n_correct / N_RELEVANT)
    return precisions, recalls
\end{lstlisting}

Scorre la lista dei risultati ordinati per punteggio. Per ogni posizione, calcola la
\textbf{precisione} (quanti dei risultati visti finora sono rilevanti) e il \textbf{richiamo}
(quale frazione delle immagini rilevanti totali è stata già trovata). La query stessa viene
saltata perché avrebbe sempre punteggio nullo e si troverebbe in cima alla classifica.
Con 100 immagini per categoria (meno la query), $N_{\text{rel}} = 99$.

\subsubsection{Funzione \texttt{interpolate\_precision}}

\begin{lstlisting}[caption={Interpolazione a 11 punti}]
def interpolate_precision(recalls, precisions):
    recalls    = np.array(recalls)
    precisions = np.array(precisions)
    interp = []
    for r in RECALL_LEVELS:
        mask = recalls >= r
        interp.append(precisions[mask].max() if mask.any() else 0.0)
    return np.array(interp)
\end{lstlisting}

Implementa l'interpolazione standard a \textbf{11 punti}: per ognuno degli 11 livelli di
richiamo $\{0.0, 0.1, \ldots, 1.0\}$, la precisione interpolata è il massimo della
precisione osservata a qualsiasi richiamo maggiore o uguale a quel livello. Questo leviga
la curva, eliminando le oscillazioni dovute all'ordine dei risultati non rilevanti.

La \textbf{Average Precision} (AP) di una query è la media delle 11 precisioni interpolate.
La \textbf{Mean Average Precision} (MAP) è la media delle AP su tutte le 1000 query ed è
la metrica principale di valutazione del sistema.

\subsubsection{Ciclo di valutazione}

Il ciclo principale carica i codebook in RAM una sola volta, poi usa ogni immagine come
query, recupera l'intera classifica (\texttt{top\_k=1000}), calcola la curva P-R interpolata
e la accumula per categoria. Al termine stampa MAP globale, MAP per categoria e salva il
grafico delle curve P-R.

% ==============================================================================
\section{Algoritmo Tree-Structured VQ (TSVQ)}
\label{sec:tsvq}
% ==============================================================================

\subsection{Motivazione}

La Flat VQ costruisce il codebook trovando $K$ cluster tutti in una volta su tutti i vettori
di una regione. Questo approccio ha due limiti:
\begin{enumerate}
    \item \textbf{Difficoltà di ottimizzazione}: trovare la migliore disposizione di $K$
          centroidi con inizializzazione casuale è un problema difficile. Con $K$ grande,
          anche ripetendo k-means 5 volte non è garantito trovare la soluzione migliore.
    \item \textbf{Ricerca lineare}: per trovare il codeword più vicino a un vettore di query
          bisogna confrontarlo con tutti i $K$ codeword, uno per uno.
\end{enumerate}
La \textbf{Tree-Structured VQ (TSVQ)} affronta entrambi i problemi dividendo il processo
in una serie di scelte binarie successive, ciascuna molto più semplice.

\subsection{Struttura ad albero binario}

Un albero TSVQ di profondità $d$ è un albero binario completo. Ogni nodo contiene un
centroide; le foglie sono i codeword finali. Con profondità $d$ si hanno $2^d$ foglie.

L'albero viene memorizzato come array piatto con \textbf{indicizzazione heap}: il nodo
all'indice $i$ ha i figli agli indici $2i+1$ (sinistro) e $2i+2$ (destro). La radice è
all'indice 0. Questo consente di rappresentare l'intero albero in un singolo array NumPy
senza puntatori. Per $d=3$: 15 nodi totali, foglie agli indici 7--14.

\subsection{Costruzione ricorsiva}

L'albero viene costruito ricorsivamente, partendo da tutti i vettori della regione:

\begin{algorithm}[H]
\caption{Costruzione TSVQ per una regione}
\begin{algorithmic}[1]
\Procedure{BuildTSVQ}{$\text{vectors},\; \text{tree},\; \text{node},\; \text{depth}$}
    \If{$\text{depth} = d$} \Return \EndIf
    \State $\text{left} \leftarrow 2 \cdot \text{node} + 1$;\quad
           $\text{right} \leftarrow 2 \cdot \text{node} + 2$
    \If{tutti i vettori sono identici}
        \State Copia il centroide del nodo corrente in entrambi i figli
        \State Continua la ricorsione sui figli con gli stessi vettori
        \Return
    \EndIf
    \State Dividi i vettori in 2 gruppi con k-means ($k=2$, \texttt{n\_init=5})
    \State Salva i 2 centroidi in \texttt{tree[left]} e \texttt{tree[right]}
    \State \Call{BuildTSVQ}{gruppo sinistro, tree, left, depth+1}
    \State \Call{BuildTSVQ}{gruppo destro, tree, right, depth+1}
\EndProcedure
\end{algorithmic}
\end{algorithm}

Ad ogni livello si divide il problema in due sotto-problemi: ognuno è un clustering binario
su un sottoinsieme di vettori già relativamente omogeneo. Un clustering binario è molto più
semplice di uno con $K$ cluster: con \texttt{n\_init=5} tentativi si trova quasi sempre
la divisione ottimale.

\subsection{Ricerca del codeword per traversata dell'albero}

Per trovare il codeword più vicino a un vettore di query, si percorre l'albero dalla radice
alla foglia:

\begin{algorithm}[H]
\caption{Traversata TSVQ}
\begin{algorithmic}[1]
\Procedure{TraverseTSVQ}{$\mathbf{q},\; \text{tree}$}
    \State $\text{node} \leftarrow 0$ \Comment{parti dalla radice}
    \For{$\ell = 1$ \textbf{to} $d$}
        \State $\text{left} \leftarrow 2 \cdot \text{node} + 1$;\quad
               $\text{right} \leftarrow 2 \cdot \text{node} + 2$
        \If{il figlio sinistro è più vicino di quello destro}
            \State $\text{node} \leftarrow \text{left}$
        \Else
            \State $\text{node} \leftarrow \text{right}$
        \EndIf
    \EndFor
    \Return $\text{tree}[\text{node}]$ \Comment{codeword foglia}
\EndProcedure
\end{algorithmic}
\end{algorithm}

Ad ogni livello si sceglie il figlio più vicino, scendendo fino alla foglia. Con $d=3$ si
eseguono solo 3 confronti binari (6 distanze calcolate) invece di confrontarsi con tutti gli
8 codeword della Flat VQ. La traversata è \textit{greedy} -- non torna indietro -- e non
garantisce di trovare il codeword globalmente più vicino, ma in pratica la perdita di
accuratezza è trascurabile.

\subsection{Confronto di complessità}

\begin{table}[H]
\centering
\begin{tabular}{lccc}
\toprule
\textbf{Metodo} & \textbf{Codeword} & \textbf{Distanze calcolate per vettore (query)} \\
\midrule
Flat VQ ($K = 8$) & 8 & 8 \\
TSVQ ($d = 3$)    & 8 & 6 \\
\bottomrule
\end{tabular}
\caption{Confronto del costo di ricerca del codeword più vicino.}
\end{table}

Il vantaggio di TSVQ cresce con il numero di codeword: a parità di $K = 2^d$ foglie, la
Flat VQ costa $K$ distanze per vettore, mentre la TSVQ ne costa $2d = 2\log_2 K$.

% ==============================================================================
\section{Implementazione TSVQ}
\label{sec:impl_tsvq}
% ==============================================================================

L'implementazione TSVQ è nel branch \texttt{TSVQ\_variant} del repository. Le modifiche
rispetto alla Flat VQ riguardano solo \texttt{build\_codebooks.py} e \texttt{retrieve.py};
preprocessing, assegnazione delle regioni, valutazione e interfaccia di \texttt{retrieve()}
rimangono invariati.

\subsection{Modifiche a \texttt{build\_codebooks.py}}

\subsubsection{Funzione \texttt{build\_tsvq\_region}}

\begin{lstlisting}[caption={Funzione ricorsiva \texttt{build\_tsvq\_region}}]
def build_tsvq_region(vectors, tree, node_idx, current_depth):
    if current_depth == DEPTH:
        return
    left_idx  = 2 * node_idx + 1
    right_idx = 2 * node_idx + 2
    if len(np.unique(vectors, axis=0)) < 2:
        tree[left_idx]  = tree[node_idx]
        tree[right_idx] = tree[node_idx]
        build_tsvq_region(vectors, tree, left_idx,  current_depth + 1)
        build_tsvq_region(vectors, tree, right_idx, current_depth + 1)
        return
    kmeans = KMeans(n_clusters=2, n_init=5)
    kmeans.fit(vectors)
    tree[left_idx]  = kmeans.cluster_centers_[0]
    tree[right_idx] = kmeans.cluster_centers_[1]
    build_tsvq_region(vectors[kmeans.labels_ == 0], tree, left_idx,  current_depth + 1)
    build_tsvq_region(vectors[kmeans.labels_ == 1], tree, right_idx, current_depth + 1)
\end{lstlisting}

Punti chiave:
\begin{itemize}
    \item \textbf{Caso base}: quando si raggiunge la profondità \texttt{DEPTH} la ricorsione
          si ferma; le foglie sono già state valorizzate.
    \item \textbf{Guardia per vettori identici}: se tutti i vettori nella partizione sono
          identici, k-means con $k=2$ non può dividere nulla. I due figli ereditano il
          centroide del nodo corrente e la ricorsione continua, producendo un albero
          strutturalmente valido anche in questo caso degenere.
    \item \textbf{Modifica in-place}: l'array \texttt{tree} viene modificato direttamente
          evitando copie intermedie.
    \item \textbf{Partizionamento tramite} \texttt{labels\_}: la selezione dei vettori per
          ciascun gruppo avviene tramite indicizzazione booleana NumPy, senza loop Python.
\end{itemize}

\subsubsection{Funzione \texttt{build\_codebooks} (TSVQ)}

\begin{lstlisting}[caption={Funzione \texttt{build\_codebooks} modificata per TSVQ}]
def build_codebooks(img):
    h, w, _ = img.shape
    n_blocks_y = h // BLOCK_SIZE
    n_blocks_x = w // BLOCK_SIZE
    features   = extract_block_features(img)
    region_ids = assign_regions(n_blocks_y, n_blocks_x).flatten()
    n_nodes = 2 ** (DEPTH + 1) - 1
    trees = np.zeros((GRID * GRID, n_nodes, 6))
    for r in range(GRID * GRID):
        region_feats = features[region_ids == r]
        trees[r, 0]  = region_feats.mean(axis=0)
        build_tsvq_region(region_feats, trees[r], 0, 0)
    return trees
\end{lstlisting}

Rispetto alla versione Flat VQ:
\begin{itemize}
    \item Il valore di ritorno è un array NumPy di forma $(16, 15, 6)$ -- 16 regioni, 15
          nodi per albero, 6 feature -- invece di una lista di matrici. Questo rende più
          efficiente sia il salvataggio che l'indicizzazione durante il retrieval.
    \item Il nodo radice (indice 0) viene inizializzato con la media di tutti i vettori
          della regione prima di avviare la ricorsione.
\end{itemize}

\subsection{Modifiche a \texttt{retrieve.py}}

\subsubsection{Funzione \texttt{region\_mse} (TSVQ)}

\begin{lstlisting}[caption={Funzione \texttt{region\_mse} TSVQ con traversata vettorizzata}]
def region_mse(query_vectors, tree):
    node_indices = np.zeros(len(query_vectors), dtype=np.intp)
    for _ in range(DEPTH):
        left  = 2 * node_indices + 1
        right = 2 * node_indices + 2
        dist_left  = np.sum((tree[left]  - query_vectors) ** 2, axis=1)
        dist_right = np.sum((tree[right] - query_vectors) ** 2, axis=1)
        node_indices = np.where(dist_left <= dist_right, left, right)
    leaf_vecs = tree[node_indices]
    return float(np.mean(np.sum((leaf_vecs - query_vectors) ** 2, axis=1)))
\end{lstlisting}

Invece di scendere l'albero un vettore alla volta, tutti i 256 vettori della regione
scendono \textbf{contemporaneamente} a ogni livello, usando operazioni NumPy vettoriali:
\begin{itemize}
    \item \texttt{node\_indices} tiene traccia del nodo corrente per ognuno dei 256 vettori.
    \item A ogni iterazione (un livello dell'albero), si calcolano gli indici dei figli e
          le distanze per tutti i vettori in parallelo, e si aggiorna \texttt{node\_indices}
          con \texttt{np.where}.
    \item L'indicizzazione \texttt{tree[left]} con un array di indici produce direttamente
          i centroidi dei figli sinistri per tutti i vettori, in un'unica operazione.
\end{itemize}
Il ciclo esegue solo \texttt{DEPTH=3} iterazioni, senza nessun loop Python sui singoli
vettori.

\subsection{Selezione della profondità}
\label{subsec:depth}

\subsubsection{Metodologia di test}

La profondità $d$ controlla il numero di codeword ($2^d$) e influenza la qualità del codebook
in due direzioni opposte: profondità bassa significa pochi codeword e scarsa capacità
discriminativa; profondità alta significa molti codeword ma foglie costruite su pochi vettori,
rendendo i centroidi inaffidabili. Si attende quindi un ottimo a profondità intermedia.

Sono state testate le profondità $d \in \{2, 3, 4, 5, 6\}$ confrontando con il baseline
Flat VQ con $K=12$ (configurazione in uso al momento di questa fase sperimentale, prima
dell'allineamento finale a $K=8$).

\subsubsection{Risultati}

\begin{table}[H]
\centering
\small
\begin{tabular}{lrrrrrrr}
\toprule
\textbf{Categoria} & \textbf{Flat VQ} & \textbf{d=2} & \textbf{d=3} & \textbf{d=4}
    & \textbf{d=5} & \textbf{d=6} \\
\midrule
dinosaurs  & 0.9096 & 0.9722 & 0.9327 & 0.8785 & 0.8278 & 0.7949 \\
horses     & 0.5214 & 0.5205 & 0.5224 & 0.5275 & 0.5185 & 0.5042 \\
buses      & 0.4693 & 0.3464 & 0.4035 & 0.4285 & 0.4343 & 0.4381 \\
africans   & 0.4685 & 0.4554 & 0.4720 & 0.4692 & 0.4713 & 0.4704 \\
food       & 0.3958 & 0.3518 & 0.3752 & 0.3972 & 0.4114 & 0.4150 \\
flowers    & 0.3600 & 0.4251 & 0.3833 & 0.3526 & 0.3296 & 0.3104 \\
elephants  & 0.3240 & 0.3454 & 0.3355 & 0.3272 & 0.3183 & 0.3096 \\
mountains  & 0.2685 & 0.2845 & 0.2835 & 0.2776 & 0.2674 & 0.2570 \\
buildings  & 0.2599 & 0.2569 & 0.2690 & 0.2678 & 0.2636 & 0.2569 \\
beaches    & 0.2263 & 0.2648 & 0.2454 & 0.2319 & 0.2196 & 0.2103 \\
\midrule
\textbf{MAP} & \textbf{0.4203} & \textbf{0.4223} & \textbf{0.4223} & 0.4158 & 0.4062 & 0.3967 \\
\bottomrule
\end{tabular}
\caption{MAP per categoria al variare della profondità TSVQ. La colonna ``Flat VQ'' usa
$K=12$ (baseline della fase sperimentale, prima dell'allineamento a $K=8$).}
\label{tab:depth}
\end{table}

\subsubsection{Analisi e scelta finale}

\begin{itemize}
    \item \textbf{d=2 e d=3 producono la stessa MAP globale (0.4223)}, entrambi superiori
          al baseline Flat VQ (0.4203). Da d=4 in poi la MAP scende monotonicamente fino a
          0.3967 a d=6, peggiore del baseline.
    \item \textbf{Categorie con distribuzione cromatica compatta} come dinosaurs, flowers,
          beaches ed elephants ottengono i risultati migliori a profondità basse. I loro
          colori sono concentrati in pochi gruppi naturali che un piccolo numero di codeword
          descrive già bene; profondità maggiori frammentano questi gruppi in modo eccessivo.
    \item \textbf{Categorie con distribuzione cromatica complessa} come buses e food
          migliorano con profondità crescente, ma non abbastanza da compensare le categorie
          che peggiorano.
    \item \textbf{Da d=5 in poi tutte le categorie peggiorano}: con $2^5=32$ foglie su 256
          vettori, ogni foglia è costruita in media su 8 vettori, troppo pochi per stimare
          un centroide affidabile.
\end{itemize}

\textbf{Scelta finale: d=3}. Pur producendo la stessa MAP di d=2, genera un codebook più
espressivo (8 foglie invece di 4), catturando più variazione intra-regione senza sacrificare
la stabilità. Con circa 32 vettori per foglia in media la stima dei centroidi rimane robusta,
e il codebook risultante è più ricco di informazione discriminativa.

% ==============================================================================
\section{Valutazione teorica dei due algoritmi}
\label{sec:valutazione}
% ==============================================================================

In questa sezione si confrontano Flat VQ e TSVQ dal punto di vista teorico. I risultati
sperimentali sono nella Sezione~\ref{sec:evaluation}.

\subsection{Complessità di costruzione del codebook}

La costruzione è una fase offline, eseguita una volta sola.

\textbf{Flat VQ} esegue k-means con $K$ cluster su 256 vettori per ogni regione. Il costo
cresce linearmente con $K$: raddoppiare i codeword raddoppia il tempo di costruzione.

\textbf{TSVQ} esegue k-means binario ($k=2$) a ogni livello. Il costo totale cresce con la
profondità $d$, ma poiché $d = \log_2 K$, il costo scala in modo logaritmico rispetto al
numero di codeword. Per $K$ grande, costruire un albero TSVQ è significativamente più
veloce che eseguire k-means con $K$ cluster in un colpo solo.

\subsection{Complessità di retrieval}

\textbf{Flat VQ}: trovare il codeword più vicino richiede confrontare il vettore di query
con tutti i $K$ codeword. Il costo scala linearmente con $K$.

\textbf{TSVQ}: la traversata richiede $d$ confronti binari, indipendentemente da quante
foglie ci siano. Poiché $d = \log_2 K$, il costo scala logaritmicamente: raddoppiare i
codeword aggiunge un solo livello all'albero, non raddoppia il costo.

\subsection{Qualità del codebook}

\textbf{Flat VQ} ottimizza globalmente la posizione di tutti i $K$ centroidi insieme. Con
$K$ grande e inizializzazione casuale è difficile trovare la soluzione ottimale: due centroidi
possono finire vicini tra loro lasciando scoperta un'area della distribuzione.

\textbf{TSVQ} ottimizza localmente ogni split binario. Dividere un insieme di vettori in due
gruppi è un problema molto più semplice: con \texttt{n\_init=5} si trova quasi sempre la
divisione migliore. I centroidi risultanti tendono a essere più bilanciati e a coprire meglio
la distribuzione.

Il compromesso è che TSVQ non garantisce l'ottimo globale: la migliore divisione binaria a
ogni livello non porta necessariamente alla migliore suddivisione complessiva in $2^d$ cluster.

\subsection{Stabilità con dati sparsi}

Con 256 vettori per regione e $K=8$ codeword, entrambi i metodi hanno in media 32 vettori
per codeword -- una stima ragionevolmente robusta in entrambi i casi.

A profondità maggiori la stabilità di TSVQ peggiora più velocemente: le prime foglie
costruite tendono a essere più popolate, mentre quelle risultanti da suddivisioni successive
possono avere pochissimi vettori, amplificando l'instabilità della stima.

\subsection{Occupazione di memoria}

Flat VQ salva solo i $K$ centroidi foglia per ogni regione. TSVQ salva anche i
$K-1$ nodi interni dell'albero, occupando quasi il doppio della memoria a parità di foglie.
Per 1000 immagini con $K=8$ e 16 regioni: Flat VQ occupa circa 6 MB, TSVQ circa 11 MB.
Entrambi i valori sono ampiamente gestibili su hardware moderno.

\subsection{Riepilogo}

\begin{table}[H]
\centering
\begin{tabular}{lcc}
\toprule
\textbf{Aspetto} & \textbf{Flat VQ} & \textbf{TSVQ} \\
\midrule
Codeword per regione       & $K$                  & $2^d$ \\
Costo costruzione          & Lineare in $K$       & Logaritmico in $K$ \\
Costo retrieval            & Lineare in $K$       & Logaritmico in $K$ \\
Qualità ottimizzazione     & Globale              & Locale binaria \\
Ricerca codeword           & Esaustiva (ottimale) & Greedy (subottimale) \\
Memoria per immagine       & $\sim R \cdot K$     & $\sim R \cdot 2K$ \\
\bottomrule
\end{tabular}
\caption{Confronto riassuntivo tra Flat VQ e TSVQ.}
\label{tab:summary}
\end{table}

% ==============================================================================
\section{Risultati sperimentali}
\label{sec:evaluation}
% ==============================================================================

\subsection{Configurazione sperimentale}

I due sistemi condividono tutti i parametri tranne la struttura del codebook:
\begin{itemize}
    \item \textbf{Flat VQ}: $K=8$ centroidi per regione, k-means con \texttt{n\_init=5}.
    \item \textbf{TSVQ}: profondità $d=3$, 8 foglie per regione, k-means binario ricorsivo
          con \texttt{n\_init=5}.
\end{itemize}
Parametri comuni: BLOCK\_SIZE=4, GRID=4, IMG\_SIZE=256, normalizzazione fissa tramite
NORM\_SCALE, spazio colore CIE L*u*v*, valutazione su 1000 query, metrica MAP con
interpolazione a 11 punti. Ogni sistema è stato valutato su \textbf{6 esecuzioni
indipendenti} per stimare la variabilità dovuta all'inizializzazione casuale di k-means.

\subsection{MAP: risultati per esecuzione}

La Tabella~\ref{tab:map_runs} riporta la MAP ottenuta in ciascuna delle 6 esecuzioni
per entrambi i metodi.

\begin{table}[H]
\centering
\begin{tabular}{lrr}
\toprule
\textbf{Esecuzione} & \textbf{Flat VQ ($K=8$)} & \textbf{TSVQ ($d=3$)} \\
\midrule
1 & 0.4198 & 0.4222 \\
2 & 0.4249 & 0.4219 \\
3 & 0.4247 & 0.4222 \\
4 & 0.4250 & 0.4222 \\
5 & 0.4252 & 0.4223 \\
6 & 0.4255 & 0.4225 \\
\midrule
\textbf{Media}  & \textbf{0.4242} & \textbf{0.4222} \\
\textbf{Min}    & 0.4198          & 0.4219 \\
\textbf{Max}    & 0.4255          & 0.4225 \\
\bottomrule
\end{tabular}
\caption{MAP per esecuzione e valori medi. Ogni esecuzione ricostruisce i codebook
da zero con una nuova inizializzazione casuale di k-means.}
\label{tab:map_runs}
\end{table}

\subsection{Curva di precisione-richiamo}

La Tabella~\ref{tab:pr} riporta la curva P-R media (su 6 esecuzioni) per entrambi i metodi.

\begin{table}[H]
\centering
\begin{tabular}{rrrr}
\toprule
\textbf{Richiamo} & \textbf{Flat VQ (media)} & \textbf{TSVQ (media)} & \textbf{$\Delta$} \\
\midrule
0.0 & 0.8690 & 0.8659 & +0.0031 \\
0.1 & 0.6342 & 0.6321 & +0.0021 \\
0.2 & 0.5362 & 0.5323 & +0.0039 \\
0.3 & 0.4758 & 0.4718 & +0.0040 \\
0.4 & 0.4261 & 0.4230 & +0.0031 \\
0.5 & 0.3842 & 0.3822 & +0.0020 \\
0.6 & 0.3447 & 0.3435 & +0.0012 \\
0.7 & 0.3081 & 0.3076 & +0.0005 \\
0.8 & 0.2726 & 0.2714 & +0.0012 \\
0.9 & 0.2366 & 0.2357 & +0.0009 \\
1.0 & 0.1786 & 0.1793 & $-$0.0007 \\
\midrule
\textbf{MAP} & \textbf{0.4242} & \textbf{0.4222} & \textbf{+0.0020} \\
\bottomrule
\end{tabular}
\caption{Curva P-R media su 6 esecuzioni. $\Delta$ = Flat VQ $-$ TSVQ (positivo indica
vantaggio Flat VQ).}
\label{tab:pr}
\end{table}

\subsection{Risultati per categoria}

La Tabella~\ref{tab:categories} riporta la MAP media per categoria su 6 esecuzioni.

\begin{table}[H]
\centering
\begin{tabular}{lrrr}
\toprule
\textbf{Categoria} & \textbf{Flat VQ (media)} & \textbf{TSVQ (media)} & \textbf{$\Delta$} \\
\midrule
dinosaurs  & 0.9333 & 0.9329 & +0.0004 \\
horses     & 0.5214 & 0.5232 & $-$0.0018 \\
buses      & 0.4582 & 0.4037 & \textbf{+0.0545} \\
africans   & 0.4668 & 0.4688 & $-$0.0020 \\
food       & 0.3858 & 0.3766 & +0.0092 \\
flowers    & 0.3785 & 0.3825 & $-$0.0040 \\
elephants  & 0.3288 & 0.3360 & $-$0.0072 \\
mountains  & 0.2722 & 0.2832 & $-$0.0110 \\
buildings  & 0.2625 & 0.2693 & $-$0.0068 \\
beaches    & 0.2345 & 0.2461 & $-$0.0116 \\
\midrule
\textbf{MAP} & \textbf{0.4242} & \textbf{0.4222} & \textbf{+0.0020} \\
\bottomrule
\end{tabular}
\caption{MAP media per categoria su 6 esecuzioni. $\Delta$ = Flat VQ $-$ TSVQ (positivo
indica vantaggio Flat VQ, negativo indica vantaggio TSVQ).}
\label{tab:categories}
\end{table}

\subsection{Analisi dei risultati}

\paragraph{MAP globale e variabilità.}
La Flat VQ ottiene una MAP media leggermente superiore (0.4242 contro 0.4222), con un margine
di soli 0.002 punti. Le due metodologie sono quindi sostanzialmente equivalenti in termini
di accuratezza globale a parità di numero di codeword.

La differenza più rilevante tra i due metodi riguarda la \textbf{stabilità tra esecuzioni}.
La TSVQ mostra un range di variazione estremamente ridotto: il risultato peggiore è 0.4219
e il migliore è 0.4225, una differenza di soli 0.0006. La Flat VQ, al contrario, presenta
un range più ampio: da 0.4198 a 0.4255, con un'escursione di 0.0057. In particolare,
la prima esecuzione di Flat VQ (MAP=0.4198) è notevolmente inferiore alle altre cinque
(0.4247--0.4255), a indicare che l'inizializzazione casuale ha prodotto almeno un codebook
di qualità scadente su alcune immagini. Questo conferma quanto atteso dalla teoria: trovare
8 cluster ottimali in un colpo solo è più sensibile all'inizializzazione rispetto a una serie
di split binari.

\paragraph{Analisi per categoria.}
Il quadro per categoria rivela una situazione più articolata rispetto alla MAP globale.
La TSVQ risulta migliore in \textbf{7 categorie su 10}: horses, africans, flowers,
elephants, mountains, buildings e beaches. I vantaggi più marcati si registrano su
mountains (+0.0110 a favore TSVQ), beaches (+0.0116) e elephants (+0.0072) -- categorie
con texture e distribuzioni cromatiche variegati che beneficiano del codebook gerarchico.

La Flat VQ prevale invece su \textbf{3 categorie}: dinosaurs (vantaggio trascurabile di
0.0004), food (+0.0092) e soprattutto \textbf{buses}, con un distacco molto marcato di
+0.0545. La categoria buses rappresenta l'eccezione principale: la distribuzione cromatica
dei bus (colori vivaci e uniformi, ma diversificati tra le immagini della categoria) sembra
beneficiare dell'ottimizzazione globale di k-means, mentre i successivi split binari della
TSVQ non riescono a catturare queste strutture altrettanto efficacemente. Questo risultato
è coerente con quanto già osservato nella fase di selezione della profondità
(Tabella~\ref{tab:depth}), dove buses era l'unica categoria a preferire chiaramente la
Flat VQ su tutti i valori di profondità testati.

\paragraph{Curva P-R.}
La curva P-R media evidenzia che il vantaggio della Flat VQ è concentrato nella zona a
\textbf{basso richiamo} (0.0--0.5), dove la precisione è più alta. In questa zona -- che
corrisponde alle prime posizioni nella classifica, le più importanti per un utente -- la
Flat VQ mantiene una precisione costantemente superiore di 0.002--0.004 punti. All'estremo
opposto, a richiamo 1.0 la TSVQ supera marginalmente la Flat VQ (0.1793 contro 0.1786),
ma le differenze in questa zona sono poco significative.

\paragraph{Conclusione.}
I due metodi si equivalgono in termini di accuratezza globale a parità di codeword ($K=8$).
La Flat VQ è leggermente superiore in media (+0.002 MAP) ma più sensibile all'inizializzazione
casuale. La TSVQ è più stabile e risulta migliore sulla maggioranza delle categorie,
ma viene penalizzata in modo significativo dalla categoria buses. Quest'ultima si conferma
un caso anomalo già identificato durante la selezione della profondità. La bilancia pende a favore della Flat VQ nella MAP complessiva.

\end{document}
